% !PRINTED 450
\weberpage{450}
%<<CONT
k\"onnen ganz von einander unabh\"angig sein; es ist aber auch der Fall nicht ausgeschlossen, dass gewisse Beziehungen zwischen ihnen festgesetzt sind, die beim Rechnen zu ber\"ucksichtigen sind.

Eine Function eines Functionenk\"orpers gilt nur dann als Null, wenn sie identisch, d. h. f\"ur alle in Betracht kommenden Werthe der Ver\"anderlichen, verschwindet.

Da wir vorl\"aufig unsere Betrachtungen nicht einschr\"anken wollen, so werden wir jetzt von \gesperrt{K\"orpern} schlechtweg sprechen, und die Objecte, mit denen die Rechnungen auszuf\"uhren sind, die sowohl Zahlen als Functionen sein k\"onnen, als \gesperrt{Gr\"ossen} oder auch als die \gesperrt{Elemente} des K\"orpers bezeichnen $^1$).

Ein \gesperrt{K\"orper} ist dann also ein System von Gr\"ossen von der Vollst\"andigkeit, dass in ihm die Gr\"ossen addirt, subtrahirt, multiplicirt und dividirt werden k\"onnen.

Wenn alle Gr\"ossen eines K\"orpers in einem zweiten K\"orper enthalten sind, so heisst der erste K\"orper ein \gesperrt{Theiler} des zweiten.

Ist $\Omega$ ein Theiler von $\Omega'$, so werden wir auch sagen, $\Omega'$ \gesperrt{ist ein K\"orper \"uber} $\Omega$.

% [sic] The registered source prints the anomalous hyphenated phrase "ein- für allemal"; preserved diplomatically.
Streng genommen, bildet die einzige Zahl ,,Null`` f\"ur sich einen K\"orper. Diesen wollen wir aber der K\"urze wegen ein- f\"ur allemal ausnehmen.

Dann ist das n\"achstliegende Beispiel eines K\"orpers der \gesperrt{Inbegriff aller rationalen Zahlen}.

Dieser K\"orper ist ein Theiler von jedem anderen K\"orper; denn jeder K\"orper enth\"alt wenigstens eine von Null verschiedene Gr\"osse $\omega$, also auch den Quotienten $\omega:\omega=1$, d. h. die Zahl 1, also auch, da alle ganzen Zahlen durch Addition und Subtraction von Einern entstehen, alle ganzen Zahlen; und aus den ganzen Zahlen kann man wieder durch Division alle Br\"uche ableiten.

Andere Beispiele von Zahlk\"orpern sind: der Inbegriff aller complexen Zahlen von der Form $x+yi$, worin $i=\sqrt{-1}$, $x$ und $y$ alle rationalen Zahlen bedeuten; ferner der Inbegriff aller (rationalen und irrationalen) \gesperrt{reellen Zahlen}, oder der Inbegriff aller \"uberhaupt existirenden complexen Zahlen $x+yi$.

Als Beispiel eines Functionenk\"orpers mag der Inbegriff aller ganzen und gebrochenen rationalen Functionen einer Ver\"ander-
%CONT>>

\par\noindent\rule{4cm}{0.4pt}\par
{\footnotesize $^1$) Vgl. des Verfassers Abhandlung: ,,Die allgemeinen Grundlagen der \gesperrt{Galois}'schen Gleichungstheorie.`` Mathematische Annalen, Bd.\,43.}
